さらに、Paroloら\cite{Par1}、Wangら\cite{Wan1}のモデルを勘案し、この指標を用いて高被引用論文を次の5クラス（FP/SB/Nm/G/O）に分類する。
うち、クラスNmについては、のちの分析のためさらにサブクラスNoを設ける。
\begin{itemize}
\item FP; Flash in the pan：Sunらは Flash in the pan の基準として、$Gs=-0.600$、$A^- = -0.300$ を示している。本研究ではこの値を参照し、$Gs < -0.600$ かつ $A^- < -0.300$ の値を持つ論文を Flash in the pan とする。
\item SB; Sleeping beauty：Sunらは $A^+$ を用いた Sleeping beauty の基準を示していないが、本研究では Flash in the pan と対照的な条件として $Gs > 0.600$ かつ $A^+ > 0.300$ の値を持つ論文を Sleeping beauty とする。
\item Nm; 成熟論文：Flash in the pan にも Sleeping beauty にも含まれない標準的かつある程度引用を得た論文の条件として、$-0.600 \leq Gs \leq 0.600$ かつ $A^- < 0$ かつ $A^+ > 0$ かつ $|A^-| > |A^+|$ を設定し、これに当てはまる論文を成熟論文とする。
\item No; 成熟論文のうちオブソレッセンスが確認されるもの：成熟論文のうちオブソレッセンスの傾向が強いものとして $A^- < -0.1$ となる論文集合を設定する。
\item G; 成長論文：被引用が成長段階にある論文として、$-0.600 \leq Gs \leq 0.600$ かつ $A^- = 0$ を設定する。
\item O; その他：以上の条件に含まれないものとして特殊な被引用形状のもの（本研究で前提としているモデルである対数正規分布に合致しないリバイバル論文等）が考えられる。
\end{itemize}
